\begingroup
\renewcommand{\arraystretch}{1.45}
\fancyhead[L]{\small HUNDERT SÄTZE ÜBER DIE NEUEN TRANSZENDENTEN.}

\Needspace{34\baselineskip}
\par\medskip\Needspace{5\baselineskip}\pdfbookmark[1]{Art. 6}{transscendentes-iv-art-6}\centerline{6.}\nobreak\smallskip
Um den Zusammenhang des Algorithmus des arithmetisch-geometrischen Mittels mit unseren Funktionen weiter zu entwickeln, bemerken wir zunächst, dass man, wenn
\[
\begin{gathered}\frac{n}{m}=\frac{(Qx)^2}{(Px)^2}\ \text{und}\ m=\mu Px^2,\\[8pt]\quad{}\qquad n=\mu Qx^2,\\[8pt]\quad{}\qquad mm-nn=\mu\mu(R)^4\end{gathered}
\]
gesetzt wird, hat


\SourceLeaf{477}

\[
\begin{gathered}
m'=\mu(Pxx)^2, n'=\mu(Qxx)^2,\\[6pt]m'm'-n'n'=\mu\mu(Rxx)^4\\[8pt]
m''=\mu(Px^4)^2,n''=\mu(Qx^4)^2,\\[6pt]m''m''-n''n''=\mu\mu(Rx^4)^4\\[8pt]
m'''=\mu(Px^8)^2,n'''=\mu(Qx^8)^2,\\[6pt]m'''m'''-n'''n'''=\mu\mu(Rx^8)^4\\[8pt]
m''''=\mu(Px^{16})^2,n''''=\mu(Qx^{16})^2,\\[6pt]m''''m''''-n''''n''''=\mu\mu(Rx^{16})^4
\end{gathered}
\]

\Needspace{18\baselineskip}
usw. oder allgemein
\[
\begin{gathered}m^{(\lambda)}=\mu(Px^{2^\lambda})^2,\\[8pt]\quad{}\qquad
n^{(\lambda)}=\mu(Qx^{2^\lambda})^2,\\[8pt]\quad{}\qquad
m^\lambda m^\lambda-n^\lambda n^\lambda=\mu\mu(Rx^{2^\lambda})^4\end{gathered}
\]

\Needspace{19\baselineskip}
und dass offenbar $\mu$ das arithmetisch-geometrische Mittel zwischen $m$ und $n$ selbst ist. Sobald wir also den Wert von $x$ bestimmen können, der vorgegebenen Werten von $m$ und $n$ entspricht, werden wir jedes Glied der Reihen
\[
\begin{array}{lll}
m,&m',&m''\mathinner{.\,.}\\[8pt]
n,&n',&n''\mathinner{.\,.}
\end{array}
\]
unmittelbar darstellen sowie die Reihen interpolieren und rückwärts fortsetzen können. Das nächstliegende Mittel zur Bestimmung von $x$ ist das folgende.


\Needspace{22\baselineskip}
Man sieht leicht, dass sich die Glieder der Reihe
\[
\begin{gathered}
h=\Bigl(\frac{Rx}{2}\Bigr)^4,\qquad
h'=\Bigl(\frac{Rxx}{2}\Bigr)^2,\qquad
h''=\frac{Rx^4}{2},\qquad
\\[8pt]\quad{}h'''=\Bigl(\frac{Rx^8}{2}\Bigr)^{\frac12},\\[8pt]
h''''=\Bigl(\frac{Rx^{16}}{2}\Bigr)^{\frac14},\qquad
h^{\mathrm{v}}=\Bigl(\frac{Rx^{32}}{2}\Bigr)^{\frac18}\mathinner{.\,.\,.}
\end{gathered}
\]
dem $x$ immer mehr nähern, sodass $x$ ihre Grenze ist. Man hat also
\[
x=h.\frac{h'}{h}.\frac{h''}{h'}.\frac{h'''}{h''}.\frac{h''''}{h'''}\mathinner{.\,.\,.}
\]

\Needspace{22\baselineskip}
Nun gilt aber nach 14
\[
\begin{gathered}\frac{h}{h'}=(Pxx)^2\ \text{und so}\ \frac{h'}{h''}=Px^4,\\[8pt]\quad{}\qquad
\frac{h''}{h'''}=(Px^8)^{\frac12},\\[8pt]\quad{}\qquad
\frac{h'''}{h''''}=(Px^{16})^{\frac14}\end{gathered}
\]

\Needspace{15\baselineskip}
Also
\[
\begin{gathered}
x=\frac{h}{(Pxx)^2.Px^4.(Px^2)^{\frac12}.(Px^{16})^{\frac14}\mathinner{.\,.}}\\[8pt]
=\frac{h}{\dfrac{m'}{\mu}.\Bigl(\dfrac{m''}{\mu}\Bigr)^{\frac12}.\Bigl(\dfrac{m'''}{\mu}\Bigr)^{\frac14}.\Bigl(\dfrac{m''''}{\mu}\Bigr)^{\frac18}\mathinner{.\,.\,.}}
\end{gathered}
\]

\SourceLeaf{478}


\Needspace{20\baselineskip}
wir haben folglich
\[
\begin{gathered}25.\quad x=\frac{mm-nn}{16m'\mu\Bigl(\dfrac{m''}{\mu}\Bigr)^{\frac12}.\Bigl(\dfrac{m'''}{\mu}\Bigr)^{\frac14}.\Bigl(\dfrac{m''''}{\mu}\Bigr)^{\frac18}\mathinner{.\,.}}\end{gathered}
\]
wofür man auch schreiben kann
\[
\begin{gathered}x=\frac{mm-nn}{16m'm''}.\Bigl(\frac{m''}{m'''}\Bigr)^{\frac12}.\Bigl(\frac{m'''}{m''''}\Bigr)^{\frac14}.\Bigl(\frac{m''''}{m^{\mathrm{v}}}\Bigr)^{\frac18}\mathinner{.\,.\,.}\end{gathered}
\]
Da sich die Glieder $m'$, $m''$, $m'''$ äußerst schnell der Gleichheit nähern, erhält man $x$ mit größter Bequemlichkeit.
\begin{center}\rule{0.25\textwidth}{0.3pt}\end{center}





\clearpage
\endgroup
